\section{Probability Space}

\subsection{Random Experiment}
\begin{itemize}
\item A \textbf{random experiment} is an experiment in which the
outcome varies in an unpredictable fashion when the experiment is
repeated under the same conditions, i.e., \emph{random}.
\item An \textbf{outcome} of a random experiment is a result that
cannot be decomposed into other results.
\end{itemize}

\subsection{Sample Space}
\begin{itemize}
\item The set of all possible outcomes for a random experiment is
called the \textbf{sample space} \(\Omega\).
\end{itemize}

\paragraph{Examples:}
\begin{enumerate}
\item Tossing a coin: \(\Omega=\{H, T\}\)
\item Tossing a coin twice: \(\Omega=\{H H, H T, T H, T T\}\)
\item Taking EEE 5544 and getting a grade:
\(
\Omega=\{A, A-, B+, B, B-, C+, C, D, F\}
\)
\item Picking an integer at random:
\(\Omega=\{\ldots,-3,-2,-1,0,1,2,3, \ldots \}\)
\item Checking the temperature (in $K$) at a random location
\(
\Omega=(0, \infty)
\)
\end{enumerate}

\begin{itemize}
\item \textbf{Finite} \(\Omega\): Examples 1, 2, 3
\item \textbf{Countably infinite} \(\Omega\): Example 4
\item \textbf{Discrete} \(\Omega\):  Examples 1--4
\item \textbf{Continuous (uncountably infinite)} \(\Omega\):  Example
5
\end{itemize}

\subsection{Event Class}
\begin{itemize}
\item An \textbf{event} is a set of outcomes (a subset of \(\Omega\))
   that we can ``quantify its happenstance'' (see below). 
\item Two special events: 
\begin{itemize}
\item \(\Omega =\)  \textbf{certain event} 
\item \(\emptyset =\) \textbf{impossible event}
\end{itemize}
\item An \textbf{event class} \(\mathcal{F}\) of a random experiment
is a collection of events.
\end{itemize}

\subsubsection{Discrete sample space}
\begin{itemize}
\item For a discrete \(\Omega\), we typically consider the event class
that is the collection of all subsets of \(\Omega\), i.e., \(\mathcal{F} = 2^{\Omega}\),
which is called the \emph{power set} of \(\Omega\).
\end{itemize}

\paragraph{Examples:}
\begin{enumerate}
\item Tossing a coin: \(\mathcal{F} =\{\emptyset,\{H\},\{T\},\{H T\}\}\).
\item Tossing a coin twice:  
\begin{equation*}
\mathcal{F} =& \left\{\emptyset,\{H H\},\{T T\},\{H T\},\{T H\},
 \{H H, T T\},\{H H, H T\}, \ldots, 
\{H H, T T, H T\}, \ldots,\{H H, T T, H T, T H\} \right\}.
\end{equation*}
\item If \( \Omega=\left\{a_{1}, a_{2}, \ldots, a_{N}\right\}\), then 
\begin{equation*}
\mathcal{F}=\left\{\emptyset,\{a_{1}\},
\ldots,\{a_{N}\}, \{a_{1}, a_{2}\},
\ldots,\{a_{1}, a_{2}, a_{3}\}, \ldots,
\{a_{1}, a_{2}, \ldots, a_{N}\}\right\}.
\end{equation*}
Altogether \(2^{N}\) events (why?).

\item For a countably infinite  \( \Omega = \left\{a_{1}, a_{2},
\ldots \right\}\), we can use the same idea as in the finite case
using the power set as the event space, i.e., 
\begin{equation*}
\mathcal{F} = 2^{\Omega} = \left\{\emptyset,\{a_{1}\},
\ldots,\{a_{N}\}, \{a_{1}, a_{2}\},
\ldots,\{a_{1}, a_{2}, a_{3}\}, \ldots,
\{a_{1}, a_{2}, \ldots, a_{n}\}, \ldots \right\}
\end{equation*}
by enumerating all subsets of \(\Omega\). However, we need a better
way to describe the power set.
\end{enumerate}

\subsubsection{Continuous sample space}
\begin{itemize}
\item For a continuous sample space, e.g., \(\Omega = (0,\infty)\),
there are uncountably infinite number of outcomes and hence the method
of enumeration used in the finite case can't be used to construct the
event class anymore.
\item Luckily mathematicians have solved the problem for us. To
discuss the solution, we need to study a bit about set operations.
\end{itemize}

\subsection{Set operations}
\begin{itemize}
  \item \textbf{Union:} \(A \cup B=\{x \in \Omega: x \in A\) or \(x \in B\}\)
  \item \textbf{Intersection:} \(A \cap B=\{x \in \Omega: x \in A\) and \(x \in B\}\)
  \item \textbf{Complement:} \(A^{c}=\{x \in \Omega: x \notin A\}\)
  \item \textbf{Difference:} \(A \setminus B=\{x \in \Omega: x \in A\) and \(x \notin B\}=A \cap B^{c}\)
\end{itemize}

\newtheorem{fact}{Fact}

\begin{fact}{De Morgan Law}
\begin{enumerate}
\item \((A \cup B)^{c}=A^{c} \cap B^{c}\).
\item \((A \cap B)^{c}=A^{c} \cup B^{c}\)
\end{enumerate}
\end{fact}
\begin{proof}
We give the proof of 1. as an example here.

\(x \in(A \cup B)^{c} \Rightarrow x \notin A \cup B \Rightarrow\left\{\begin{array}{l}x \notin A \\ x \notin B\end{array} \Rightarrow\left\{\begin{array}{l}x \in A^{c} \\ x \in B^{c}\end{array}\right.\right.\)\\
Thus \((A \cup B)^{C} \subset A^{C} \cap B^{C}\)

\[
\Rightarrow x \in A^{c} \cap 1
\]

\(x \in A^{C} \cap B^{C} \Rightarrow\left\{\begin{array}{l}x \in A^{C} \\ x \in B^{C}\end{array} \Rightarrow\left\{\begin{array}{l}x \notin A \\ x \notin B\end{array} \Rightarrow x \notin A \cup B \Rightarrow x \in(A \cup B)^{C}\right.\right.\) Thus \(A^{C} \cap B^{C} \subset(A \cup B)^{C}\)

Therefore \((A \cup B)^{C}=A^{C} \cap B^{C}\).
\end{proof}
\begin{fact}
Union and complement are enough to specify all other set operations.
\end{fact}



\begin{itemize}
  \item[-] Need a rath object to allow description of a random experiment.
  \item[-] Probability space \((\Omega, F, P)\) Sample space
  \item[] probability measure
  \item[-] Sample space \(\Omega\) : set that contains all possible ants of the random experiment.
\end{itemize}
